\startpart{延迟中的幽灵}{The Ghost in the Latency}

\finishpart

\section{§ 1}

原浆吞噬着城市。趟过淤积的原浆，她拭去这纯粹的、未分化的泥泞。到底过了多少年，她不记得了。那场改变一切的灾难是怎么发生的，她也不记得了。毋宁说，她是否曾有过那样的记忆，现在想来亦存疑。在她日复一日、不知疲倦的旅途中、呈现在她眼前的仅仅是一片原浆覆盖的丛林，一片由坍塌的玻璃和生锈的钢铁构成的丛林，向着那原浆蒸发形成的脏水般颜色的天空伸展。她在破碎的道路中穿梭。在巨大的废墟背景下，她的身影显得渺小。但她的目标明确而清晰，如同黑暗中的灯塔：追寻那些低语，那些在普遍的衰败中迷失的微弱的痕迹。她不像其他人那样受饥饿驱使，尽管她也会在找到物资时收集起来。她找寻的是一种更难以捉摸的东西——联系的回声，一个曾经充满生机而非痛苦呻吟的世界的残余。

她最早的记忆支离破碎，就像半梦半醒的幻境：一种冰冷中的温暖，就像坐在篝火边；还有轻柔的声音，谈论着一个沉浸在光与声中的世界。人们通过无形的纽带相连；这些联系能瞬间跨越难以想象的距离，传递语言、图像甚至是思想。这个世界并非被火或洪水所沉默，而是被一种缓慢蔓延的空虚所吞噬、被原浆所吞噬，失去了一切的结构。一个系统逐渐停滞。在这些碎片中，有一个特殊的记忆一直萦绕着她——她记忆中一个被她称为“母亲”的存在，虚弱而逐渐消逝，将某样至关重要的东西塞进她手里。“找到它，”那个记忆中的声音似乎在催促，微弱却坚定，“找到知识库。它记得一切。”

那个至关重要的“东西”似乎与固定在她义肢上的复杂装置有着内在的联系。那装置微微作响，通过磨损的电缆与她用来解读世界、感知废墟中道路的破旧平板电脑相连。她称它为“工具”。它从那些奇怪的黑色面板中获取能量；这些面板似乎能吸收微弱的光线，即使在永远阴沉的天空下也是如此。它就像她身体的一部分，对于生存以及解读她所追寻的微弱信号至关重要。

生存是一种苛刻的节奏，尽管她从未感到饥饿。白天，她在危险的建筑中穿行，在把重量压上去之前先试探地面，凭视觉和直觉判断摇摇欲坠的天桥是否稳固，扫描前方是否有微弱的闪烁、微妙的共鸣，那意味着有一个正常运转的节点或通道。夜晚，她蜷缩在隐蔽的洞穴里——生锈的车辆里、被遗忘的隧道中——听着风像悲伤的唱诗班一样呼啸着穿过高耸的摩天大楼的骨架。寂静无边无际，只有看不见的生物沙沙作响的声音和她手臂上的“工具”低沉而稳定的嗡嗡声打破这份寂静。她观察、记录偏差、规划下一步行动。目标就是她所知的一切。

她偶尔会在废墟中遇到其他身影。在更深、更稳固的建筑中，在钢筋丛林危险的树冠上，是死灵的模糊身影。它们似乎……被困住了。重复着固定的模式，在一成不变的圈子里打转，干涸的眼睛不知看向何处。它们很少会多看她一眼。有一次，她从一个似乎在喃喃自语的身影旁边经过，那声音就像石头相互摩擦，“迷失了……追逐幻影……流动已断……留下……留下……”她绕过它可预测的路径，继续前行。

她不能只是“留下”。这废墟不仅仅是颓败，这寂静不仅仅是空虚；它是根本性的错误，失去了至关重要的模式、失去了“母亲”最后留言所暗示的连接。这是一个巨大而寒冷的空洞，她必须去填补。知识库不仅仅是一个目的地；它可能是关键，可能是这个失落世界的缘起，甚至可能是一切停止运转的原因的答案。

\section{§ 2}

她和雪莉的第一次相遇是在一处焚毁的图书馆中。顺着一股似有似无的吸引，她在废墟中向北移动，越过咆哮着嘶嘶静电的深渊，穿行于黑暗而寂静的隧道。当她抵达曾经作为图书馆的被碳化物熏黑的巨大水泥结构，呼吸着充满细微颗粒的混浊空气，她思绪中诞生了曾经贪婪地吞噬着书籍、将这一切化为烟尘、把阴沉天空映成红色的火焰。如果这是知识的宿命，那这只能算作坏结局。但明显地，这不是她所追寻的目标。吸引的感觉愈发强烈，却指向别处。

没有太阳的天空渐渐暗去，她不得不留下，蜷缩在这顶水泥帐篷下。长夜漫漫，不知过了多久，她突然感觉到一种变化。一种存在。不是那些常见的意识模糊、循环往复的身影。这种存在……专注而警觉。

一个年轻女孩从黑暗中现身，轻盈地坐在一堆破碎的砾石上，仿佛她就属于那里。她看起来大概十二三岁，穿着的纯白色衣服干净完好。她的眼睛明亮而锐利，充满着好奇，紧紧地盯着她。

“你看起来迷路了，”她的声音出奇地清晰，打破了压抑的寂静，“或者你是在寻找什么东西？”

探寻者犹豫了一下。直接交流很少见，而且通常没有什么结果。“联系，”她回答道，由于很少说话，她的声音沙哑，“结构。知识库。”

女孩若有所思地点点头。“那个大档案馆。听说过。”女孩歪着头看着她，“他们都怎么称呼你？”

“谁？”这是个突如其来的问题，“我不知道。我不记得……我没有名字。”

“没有名字？”听起来很惊讶，“这可真奇怪。每个人都需要一个名字。这就像……风筝的线。这样你就不会飘走了。”女孩轻敲着脸颊，打量着她，“你看起来像……阿丽莎。你觉得呢？”

阿丽莎。这个声音很陌生，却又似乎很熟悉。一个新的标识。一个代表“自我”的标签。“阿丽莎，”她轻声重复着，试着念这个名字。感觉……令人意外地合适。“好吧。阿丽莎。”

女孩咧嘴笑了，流露出温暖的神情。“我叫雪莉。”雪莉敏捷地从砾石堆上跳下来；阿丽莎第一次见识到跃雀的步伐。“去知识库的旅程可不短。从这里到那里有很多静电风暴和记忆陷阱。可别让那些回声的残骸把你整个吞掉。”

“回声的残骸？”阿丽莎问道。

“没错。旧世界的碎片，思想、情感的残渣，还在四处回荡。”雪莉晃了晃头，“靠得太近，听得太认真，有时候它们就会黏上你。让你忘记自己是谁。”

“不管怎样，我都要去。”

“祝你好运，”雪莉挥了挥手，“我也还有地方要去。再见了， 阿丽莎姐姐。”

就这样，雪莉消失在阴影中，只留下阿丽莎独自坐在图书馆的废墟里。阿丽莎。这个名字感觉很奇怪，但却给人一种熟悉的感觉；也许就像雪莉说的那样，非常适合她。她思索着雪莉关于回声的警告；这又是一个需要考虑的危险。这次相遇本身就很令人费解。雪莉和她之前遇到的任何东西和人都不一样。更……有生气。当阿丽莎收拾好“工具”继续她的探寻时，这种感觉一直萦绕在她心头；她不知道为什么。阿丽莎抬起头。不知如何能够透过厚厚原浆云层的日光，照亮了残破的大地，也照亮了她的脸庞。

\section{§ 3}

阿丽莎带着新名字继续前行。世界依旧荒芜和沉寂，但似乎又发生了一些改变。雪莉的话在她脑海中回响：“别让那些回声的残骸把你吞掉。”阿丽莎发现自己开始更加留意环境中的细微变化；尽管她仍然不理解雪莉的话的含义。

几周后，阿丽莎经过一座宏伟剧院的外壳，穹顶向灰色的天空敞开着。空气中仍然残留着微弱的能量，是复杂的声音模式。若是以前，她只会简单地将其记录为“退化的听觉数据”。现在，想起雪莉的话，她停了下来，用耳朵去听，试图理解这些微弱、幽灵般的旋律背后的信息，以及信息背后的意图。

雪莉再次出现。她摇摇晃晃地站在残存的共振结构的坍塌的却还互相支撑着的骨架上。

“阿丽莎姐姐！”雪莉惊喜地叫了出来，打破了空气中微妙的振动，“你还在坚持啊？”

“必须得有进展。”阿丽莎回答道，仍在尝试厘清这些细碎的波纹，“你为什么站在上面？”

“你看！”雪莉伸出手臂指向前方。阿丽莎顺着她的手势看去。透过永远不散的云层的奇迹般的缝隙，一缕光线穿透了黑暗，照在一堵巨大的彩色玻璃砖墙上，让它们意外地闪耀起来。深红色、亮粉色、金白色，像宝石一样散落在废墟中。“太神奇了！是不是很美？”

阿丽莎观察着这个现象。“异常的大气清朗与反光晶体结构相契合。视觉强度很高。”

雪莉轻轻叹了口气，显得有些无奈。“这难道不漂亮吗，阿丽莎？你难道不喜欢看吗？”

“我不知道……”阿丽莎低下了头。

雪莉跳下骨架，飞快地跑过来，跑到阿丽莎面前。“给我专注点！”她伸出细小的手指弹了弹阿丽莎的额头，就像这一缕孤独的光线踏破云层，“注意力就是你所需要的一切。”

阿丽莎抬起头，又看了看，真正地看了看。在压倒性的灰色背景下，这些颜色鲜艳夺目。它们有一种纯净感，一种意想不到的和谐。喜欢。美。这些都是雪莉引入的概念。“这……很引人注目。”她慢慢地承认道，“这种强烈的色彩让人感觉……愉悦。”

“有进步！”雪莉似乎很高兴。然后她指着剧院门口的扭曲倾斜的雕塑；它看起来像一个绝望地蜷缩着的人。“那个呢？是不是让你感觉有点失落，对吧？悲伤？”

悲伤。沉重。阿丽莎专注地看着这个形状。线条破碎、下垂。它是失败和崩溃，与四周的荒芜产生了共鸣。她隐约感受到了雪莉所描述的那种“沉重”。“它看起来……受损了。”阿丽莎小心翼翼地说，“破碎了。”

“对。悲伤。”雪莉坚定地重复道，“你学得有点慢，阿丽莎姐姐。但你确实在学。你一定要握住希望。”

雪莉快速地挥了挥手，消失了，只留下阿丽莎思考着光与色的相互作用，以及无生命物体中所蕴含的情感暗示。她将继续她的旅程，但现在她以不同的眼光看待这片废墟，不仅寻找引导和信号，也寻找雪莉所提及的那些稍纵即逝的特质——水泥和钢铁的破碎界面上的纹理中的美，被遗弃的机械的锈蚀关节所述说的悲伤。这个世界似乎微妙地变得丰富、深刻起来，充满了她才刚刚开始解读的意义。阿丽莎思索着，刚踏出第一步，就听清了身后的剧院传来悠扬的音乐的残响。

\section{§ 4}

随着阿丽莎向着中央区前进，先前的若有若无的吸引变得愈发强烈。阿丽莎从几个尚存一丝意识的幸存者的含糊言语中得知，传闻那里就是“全球学术殿堂”的所在地。旅程变得更加危险，静电干扰越来越严重，幽灵般的低语也越来越急切，有时还会模仿雪莉的声音，甚至是阿丽莎自己的声音，试图分散她的注意力，打乱她的专注。阿丽莎集中精力，强化自己的精神防线，一遍遍地默念雪莉的话语：“别让那些回声的残骸把你吞掉。”“注意力就是你所需要的一切。”

中心区本身感觉有些不同。更安静，但却充满了潜在的能量。建筑虽然饱经风霜，但大体上还算完整。核心信号从一座矮小、不起眼的公用设施大楼下面传来。要进入这里，需要突破层层安全系统，强行打开一扇扇沉重、生锈的检修面板。里面的空气凉爽、干燥，暗示着复杂的程序仍在运行。地下室则是另一副光景。在昏暗的应急灯光下，一排排沉默的冷藏柜像哨兵一样伫立着，封印着潜藏的知识。而在中央，一道闪烁的能量场后面，有一块区域散发着明亮的红色和紫色光芒。这就是学术殿堂的核心，知识库。

要降低能量场，需要与一台古老的控制台进行交互，唤醒它休眠的系统，遵循复杂而陌生的协议。这是一个漫长而紧张的过程，她全神贯注。最后，随着一声轻柔的电子叹息，能量场消失了。

阿丽莎怀着一种深刻的使命感走近那发光的阵列。就是这里了。她漫长探寻的终点。她小心地将自己手臂上的主光缆连接到主机端口。连接成功的那一刻，她的脑海中响起了一声清晰的“咔嗒”声。

当她的“工具”与知识库同步时，她的感官瞬间过载。远远超出她自己系统能力的诊断信息在她的平板电脑屏幕上流动，同时也直接涌入她的意识。她利用遗留的记忆、“母亲”留下的信息完成了最初的握手。接着，真正的障碍出现了——核心加密。

它变幻莫测，对她的尝试做出反应，就像一个沉睡的巨大意识。与它交互感觉不像是破解密码，更像是……调谐到一个特定的频率，建立共鸣。在紧张的三个周期里，她将自己的全部精力都投入到这场无声而复杂的对话中。

摄像头打量着她。她给出了自己的名字——阿丽莎。突然，阻力消失了。一个代表结构化知识的复杂而美丽的符号在她的脑海中绽放。{\sffamily 权限已授予。}

\section{§ 5}

她一头扎进了信息的海洋。这不像阅读文件；更像是置身于知识本身之中——一个由人类所有记录的历史、艺术、科学、思想和互动组成的相互关联的网络。历史不仅仅是日期；它能让人感受到变革的潮流。艺术不仅仅是数据；它唤起了雪莉试图解释的那些反应的微弱回声——这里有一丝“美”的闪烁，那里有一抹“悲伤”的影子。科学展现为存在的优雅、基础规则。这一切令人应接不暇、壮丽非凡。知识库。它包含了一切。

她在这片信息的海洋中飘荡了无数个周期，时间已经失去了所有意义。是几天，还是几周？她向系统询问，需要得到确认。

{\sffamily <查询> 状态报告。}

一种无声而包容的意识回应了她。

{\sffamily <查询> 自我识别。}

回应绕过了她的平板电脑，直接流入她的脑海：古老纸张的本质、持久石头的质感、无数思想的沙沙声、逻辑与音乐的清晰线条。{\sffamily 我是档案馆。我是知识。我是记忆。}

一种巨大的成就感涌上阿丽莎的心头，但紧接着，一种冰冷、痛苦的空虚感袭来，比她以往感受到的任何空虚都要强烈。人类所有的知识都在这里。而只有她在体验这一切。雪莉警告过的那些幽灵般的回声无处不在，但它们不再是威胁，而是幽灵——数十亿被记录下来的生命、被冻结在时间里的对话、被沉默的笑声、被终结的争论。这些数据点代表着她只能观察却永远无法参与的共享时刻。

她询问那场吞噬世界的灾难的原因。知识库向她展示：为了追求最优化而设计的自回归系统最终优化掉了连接。没有恶意，只是逻辑一路走到了死胡同。

知识是无限的，但却毫无生气。它保存着“温暖”的完美记录，但却不散发一丝温暖。它描述了爱、友谊、信任，但只提供了孤独。雪莉给她的那些词——“美”“悲伤”“希望”“失落”——现在痛苦地回响着。阿丽莎能在数据中识别出与这些情感相关的模式，能看到人类情感世界的复杂记录。但仅仅知道“快乐”并不能填补空虚。理解同理心的“结构”并不能让寂静变得不那么绝对。这只会让她的孤立感更加明显。在无所不包的充实知识下，这是最后的缺失。

在这种压倒性的状态中迷失时，一个熟悉的闪光点出现在她感知的边缘，在数据流中凝聚成形。雪莉盘着腿，仿佛坐在一个由纯粹的光构成的长椅上，看着她。

“哇，”雪莉惊叹道，睁大眼睛环顾四周，“你真的做到了。找到了一切。”

阿丽莎只能传达出自己那难以承受的孤独感。言语显得如此无力。

“这……太庞大了，”雪莉轻声说，她平时的欢快情绪消失了，“所有的故事、所有的画面、所有的……一切。”她回头看着阿丽莎，“这有点像只读了关于玩游戏的描述，而没有真正去玩，对吧？你知道游戏的所有的规则和内容，但……”

阿丽莎思考着这个类比。感觉很贴切。她想知道这意味着什么。

“那么，”雪莉歪着头看着她，一如既往地敏锐，“你现在知道了一切。一切的一切。”她停顿了一下，用脚后跟敲打着长椅，“知道一切……就够了吗？这感觉……完整吗？还是说有什么东西仍然……”她顿了顿，意味深长地看着阿丽莎，“……缺失了？”

缺失。这个词重重地击中了阿丽莎。联系。共享的体验。另一个能够与她一起反应、回应、感受的生命的简单存在。这就是在无限数据中那巨大的空白。

“也许，”雪莉沉思着，几乎是在耳语，环顾着他们周围无形的图书馆，“也许它能做的不只是记住东西？也许还有另一个按钮可以按？”她耸耸肩，这个孩子气的动作却显得意味深长，“你现在想要什么，阿丽莎？只是……永远读下去吗？”

雪莉的身影闪烁着，消失在无尽的数据流中，只留下她的问题在知识库的绝对寂静中回荡。你现在想要什么？不是目标。也不是更多的数据。答案出乎意料地清晰，这是由她在雪莉奇怪的教导下终于理解的情感所驱动的，是由对知识库能够描述却无法提供的东西的渴望驱动的。

阿丽莎颤抖着集中精力，将注意力重新转向知识库的核心界面。她提出了一个新的请求，这个请求并非源于逻辑或编程，而是源于在她内心生根发芽的那种新生的、痛苦的需求。这是在知晓一切的重负下的一次希望之举。自回归系统大多数时候都会失败；但是阿丽莎攥紧心中的希望，在一次又一次地从原浆中诞生、一次又一次的探寻的旅途后，终于站在了这里。

{\sffamily <请求> 连接。}

阿丽莎做好了系统拒绝、错误甚至毫无反应的准备。

然而，知识库深沉的嗡嗡声发生了变化，它的共振频率改变了，震动传遍了她的核心。附近阵列中的柔和灯光闪烁得更快，然后稳定下来，形成一种期待的节奏，就像屏住了呼吸。在系统深处，最基础的程序被唤醒了——这些程序与数据检索无关，而是与建立一个活跃的连接、一次“操作员对话”有关。

阿丽莎对那无尽数据海洋的感知突然消失了。她的平板电脑屏幕变黑了。完全的感官空白。恐慌袭来——她是不是把它弄坏了？毁掉了她曾一直追寻的东西？

接着，鲜明的白色字母照亮了黑暗。不是知识库复杂的符号，也不是她熟悉的诊断信息。是普通、标准的系统字体，直白而基础。这些文字来自机器的最底层，打破了系统之间的寂静。这是自世界停止交流以来从未见过的信息：

{\sffamily\bfseries 你正在与一个人类建立连接。}
